% section:abstract
We examine the conditions under which a common graded state space can
support particle and geometric readouts. The Quantum Sphaera construction
is formulated here as a set of conditional interfaces rather than a
completed identification with physical fields. For a real orthogonal
involution, the associated compensator gives an exact projection, but
does not determine a three-generation mass spectrum. We derive the
additional equal-norm condition for the Koide root ratio and express it
as an orthogonality condition in a declared HC representation. Direct
odd-input and cyclic-equivariant readouts exhibit obstructions, while
positive quadratic readouts require a restricted state domain. A
projection of the dynamics supplies a separate autonomy condition for
observable variables. Together these results isolate the operator,
metric and state-selection data needed to turn an algebraic framework
into a physical model. No lepton masses, absolute mass scale or
gravitational field equations are predicted by the present construction.
% section:question
Attempts to describe different physical sectors on a common mathematical
carrier face a recurring difficulty: a shared state space does not by
itself specify the observables that distinguish those sectors. The same
algebra can support different representations, couplings and measurement
maps. A physical interpretation therefore requires more than a matching
number of components or a symmetry with the expected name.

The Quantum Sphaera Companion develops a common-carrier proposal using
biquaternionic and division-algebra structures, an involution and a
compensating projection \cite{QSC411}. This paper extracts a connected
question from that broader programme: under what conditions can the
graded structure be transferred to particle, mass and geometric
readouts without inserting the desired observables into the transfer?
The emphasis is on explicit conditions and consistency constraints.

The mass sector makes this distinction particularly sharp. A relation
among three positive roots can be represented algebraically in many ways,
but a representation selected after the roots are specified is not a
mechanism for their values. We use the Koide ratio as a test of this
interface, derive its missing norm condition, and compare linear and
quadratic readouts. The same distinction reappears in projected dynamics:
an equation on the full carrier need not induce an autonomous equation
for the variables identified as observables.

The scope is consequently conditional. The algebraic projection and
the stated comparison results are available; the identification of a
native carrier, its couplings and its physical observables remains to be
completed. This distinction also governs the particle assignments and
geometric constructions discussed below.
% section:core
Let $H$ be a real Hilbert space and let $J:H\to H$ be an orthogonal
involution. Since $J^{-1}=J=J^*$, the operators
\begin{equation}
 C=\frac{I+J}{2},\qquad P_-=\frac{I-J}{2}
\end{equation}
are complementary orthogonal projections. Indeed,
$C^2=(I+2J+J^2)/4=C$, and similarly $P_-^2=P_-$, while
$CP_-=0$ and $C+P_-=I$. We refer to $C$ as the HC compensator.
It selects the $J$-even part of a state; $P_-$ selects the odd part.

This grading is the minimal structure needed for the subsequent tests.
It does not identify either range with a particle species. In the
Companion, a biquaternionic carrier and extensions involving
$\mathbb C\otimes\mathbb H\otimes\mathbb O$ motivate proposed
particle and frame assignments. Division algebras provide a mathematical
setting for such representation questions \cite{Baez2002}. A positive
Hilbert metric used in a readout calculation must nevertheless be
distinguished from a Lorentzian spacetime metric.

The scalar field also matters. An antiunitary involution on a complex
Hilbert space is generally only real-linear when averaged with the
identity. Here $J$ is explicitly real orthogonal; when complex amplitudes
are used, we mean its declared complex-linear extension. No equivalence
with an inherited antiunitary convention is presumed. Likewise, a change
between the proposed matter, light and tachyon frames requires an
explicit map, not merely a relabelling of components.
% section:method
The Companion uses QAM as an object and assertion language, MML for
typed instructions and MMA for their execution semantics. For the
physical arguments in this paper, their essential role is to keep an
assertion attached to its domain, assumptions, calculation and remaining
interpretation requirements. A guard denotes a condition still needed
for a stated inference. A finite calculation and a general proof have
different scopes, even when both concern the same formula.

The projection identity provides a complete elementary example. Its
inputs are a Hilbert space, $J^2=I$, $J^*=J$ and the definition of $C$;
the calculation in the preceding section proves $C^2=C$. An exact
finite instance is
\begin{equation}
 P=\frac12\begin{pmatrix}1&1\\1&1\end{pmatrix},\quad
 M=\frac14\begin{pmatrix}5&3\\3&5\end{pmatrix},\quad
 N=\begin{pmatrix}1&0\\0&2\end{pmatrix}.
\end{equation}
Direct multiplication gives $P^2=P$, $[M,P]=0$, and
\begin{equation}
 [N,P]=\frac12\begin{pmatrix}0&-1\\1&0\end{pmatrix}\ne0.
\end{equation}
An exact rational check therefore confirms the projection and one
commuting transport while rejecting the other transport as a commuting
example. This rejection is a control result, not a failure of the
projection theorem. None of these matrix equations identifies a native
field, a physical transport or a mass operator.

Ordinary mathematical notation is retained throughout the main text.
The accompanying assertion records specify precise source versions,
prerequisites and finite checks; they allow the chain of reasoning to
be inspected without treating content identity as proof. The original
mass-readout arguments are collected in Appendix~A, while the separate
audit records retain the detailed provenance and unresolved obligations.
% section:particles
The algebraic carrier suggests several representation-theoretic
ingredients: unit quaternions provide an $SU(2)$ action, complex scalars
provide phase transformations, and the octonionic setting motivates
colour representations. These ingredients are useful starting points,
but they do not uniquely fix the physical gauge theory. In particular,
a common complex phase is not by itself the electromagnetic charge
operator or a derivation of hypercharge.

The distinction between colour and generation is equally important.
A three-dimensional subspace carrying a colour action cannot be
identified with three generations merely by counting its basis vectors.
The Companion's inherited slot assignments require additional choices
of embeddings, charge operators and the relation between internal
representations and physical states. We therefore retain the gauge and
generation discussion as a conditional assignment problem, not as a
derivation of the complete Standard Model.

A more explicit comparison construction starts from a seed isometry
$T_0:\mathbb C^3\to\mathcal K$ and an orthogonal closure projector
$P_c$ on a declared common carrier. If $T_0^*P_cT_0$ is positive
definite, then
\begin{equation}
 S=P_cT_0(T_0^*P_cT_0)^{-1/2},\qquad S^*S=I_3.
\end{equation}
This constructs three orthonormal closed states under its hypotheses.
It does not determine the seed embedding or explain why those states
are lepton generations. We distinguish $P_c$ from the HC projector $C$:
closure within a selected space and even/odd HC grading are separate
structures. This distinction becomes necessary when a coupling is
used to define mass readout.
% section:masses
\paragraph{A necessary norm relation.}
Let $x_i=\sqrt{m_i}\geq0$, with $x\ne0$, and define
\begin{equation}
 K=\frac{\sum_i x_i^2}{(\sum_i x_i)^2},\qquad
 A=\frac{x_0+x_1+x_2}{3},\qquad x=A(1,1,1)+d.
\end{equation}
Since $\sum_i d_i=0$, the constant part and $d$ are orthogonal, so
\begin{equation}
 K=\frac13+\frac{\|d\|^2}{9A^2},\qquad
 K=\frac23\ \Longleftrightarrow\ \|d\|^2=3A^2.
\end{equation}
Thus the Koide value requires equality of the symmetric and traceless
squared norms. Cyclic symmetry preserves these norms separately and
does not fix their ratio. For example, $x=(1,2,3)$ gives $K=7/18$,
while equal roots give $K=1/3$.

\paragraph{A conditional HC representation.}
On Euclidean three-root space set $e=(1,1,1)/\sqrt3$ and let $J_m$
be an orthogonal involution. For $C_m=(I+J_m)/2$, suppose
$x=\alpha C_me$ is nonzero and componentwise nonnegative, with
$\alpha>0$. Writing $h=\langle e,J_me\rangle$ gives
\begin{equation}
 \|C_me\|^2=\frac{1+h}{2},\qquad
 K=\frac{2}{3(1+h)}.
\end{equation}
Consequently $h=0$ is the additional condition for Koide. The projector
axioms permit $J_m=I$, for which $h=1$ and $K=1/3$; they do not select
the desired value. Choosing $J_m$ after specifying the target roots
cannot supply the missing physical selection rule.

An explicit route from the original carrier to this representation
would require a bounded coisometry $B:H\to\mathbb R^3$, $BB^*=I_3$,
whose kernel is $J$-invariant. Then $J_m=BJB^*$ is an orthogonal
involution and $BC=C_mB$. If an independently supplied seed satisfies
$Bf_0=e$ and the physical root law is $x=\alpha BCf_0$, the criterion
becomes
\begin{equation}
 \langle B^*e,JB^*e\rangle=0.
\end{equation}
Neither the name of a readout nor the number of its components supplies
$B$, its metric control, this seed law or the orthogonality selection.
Moreover, a purely odd input obeys $Cf=0$ and cannot yield nonzero roots
through this direct compensated linear route.

\paragraph{Positive nonlinear readout is a separate requirement.}
For a nonzero pair of complex amplitudes $\chi=(u,v)$ define
$X_0=|u|^2+|v|^2$, $X_1=|u|^2-|v|^2$,
$X_2=2\operatorname{Re}(\bar uv)$ and
$X_3=2\operatorname{Im}(\bar uv)$. Then
$X_0^2=X_1^2+X_2^2+X_3^2$. With the three-cycle $U$, an illustrative
Hermitian root operator is
\begin{equation}
 L_\eta=X_0\frac{I_3}{\sqrt3}
 +\eta X_1\frac{U+U^*}{\sqrt6}
 +\eta X_2\frac{i(U-U^*)}{\sqrt6},\qquad \eta\geq0.
\end{equation}
Only where $L_\eta\geq0$ may its eigenvalues be read as principal
roots. There
\begin{equation}
 K=\frac{1+\eta^2[1-(X_3/X_0)^2]}3.
\end{equation}
Positivity on every input requires $\eta\leq1/\sqrt2$ and gives
$K\leq1/2$. The isometric coefficient $\eta=1$ can reach $2/3$
on restricted positive phase chambers, but the preparation and
preservation of such a chamber is additional physical information.
Squaring a signed $L_\eta$ does not remove this issue: its physical
principal root becomes $|L_\eta|$, changing the root sum.

The recent HC-balance comparison illustrates the transfer problem
more directly. If the two-mode grading is $J=\operatorname{diag}(1,-1)$,
the balanced state $(1,i)/\sqrt2$ gives $X_1=X_2=0$ and hence $K=1/3$
for this readout. More generally, an antipodal pair of root operators
$cI_3\pm D$, with $c>0$ and $\operatorname{Tr}D=0$, cannot both
be positive Koide readouts. Both positive implies eigenvalues of $D$
in $[-c,c]$ and $\operatorname{Tr}D^2\leq2c^2$, so both ratios
are at most $5/9$. This is a restriction on this isotropic readout
pair, not on arbitrary mass mechanisms. It shows why state balance
and the positive three-generation transfer must be derived together.
% section:geometry
The same interface issue appears in the dynamics. Suppose a finite
realification of the full state satisfies $\dot u=Gu$ and let
$r=T(t)u$ be a surjective real-linear extraction. For a right inverse
$R_s$, write $u=R_sr+z$, with $Tz=0$. Differentiation gives
\begin{equation}
 \dot r=(\dot T+TG)R_sr+(\dot T+TG)z.
\end{equation}
The observable equation is independent of the discarded lift precisely
when
\begin{equation}
 (\dot T+TG)\ker T=0.
\end{equation}
This condition is additional to an algebraic projector identity. In its
absence, two full states with the same readout can have different
readout derivatives. Eliminating the hidden variables may then produce
memory terms rather than an autonomous local equation.

The Companion's geometric programme studies induced coframes, local
field comparisons and compatible constraint data. Finite normal-source
calculations can test whether prescribed tangential data determine a
normal increment. Such a calculation is not yet a regional existence
theorem, a derivation of the native action or an identification of an
internal positive comparison metric with a spacetime metric. The
reconstruction must respect common domains, conservation conditions
and the ordering of the original operators.

This identifies a concrete link between the mass and geometry branches:
both require a represented extraction with controlled discarded degrees
of freedom. A formal master equation supplies an organizing equation;
well-posed evolution and the physical meaning of its projected fields
are further obligations.
% section:limits
The results above locate the main missing information in the transfer
from a common carrier to observables. An HC projector gives an exact
grading. Koide additionally requires a metric-preserving root
representation and a selected norm relation. Positive nonlinear
readouts introduce state-domain restrictions, while geometric
projection introduces an independent autonomy condition. None of
these requirements is fixed by the dimensionality of a representation.

A recent comparison route couples this question to an arithmetically
defined operator. For $\sigma>1$, the logarithmic derivative of zeta
defines the real odd function
\begin{equation}
 a(t)=\operatorname{Im}\frac{\zeta'(\sigma+it)}{\zeta(\sigma+it)}
 =\sum_{n\geq2}\Lambda(n)n^{-\sigma}\sin(t\log n),
\end{equation}
where $\Lambda$ is the von Mangoldt function \cite{DLMF274}.
Multiplication by $a$ anticommutes with parity $R$. Thus the
selfadjoint operator $T_a=-iRA_a$ has symmetric positive and negative
spectral subspaces. A state restricted to either nonzero spectral side
is parity-balanced, including when the spectrum is continuous.
The construction $f(t)=Na(t)^2e^{-t^2}$,
$\psi=(f+i\operatorname{sgn}(a)f)/\sqrt2$ gives a regular normalized
example on the positive side. Its parameter and Fourier variances
are finite. The corresponding Robertson relation is compatible with
the balance, but need not be saturated \cite{Robertson1929}.

This example establishes mathematical compatibility, not physical
selection. The positive spectral restriction, its preparation, the
window and the interpretation of $T_a$ as energy are additional
choices; the parameter is not identified with physical time or position.
The convergent prime series also does not justify substituting
$\sigma=1/2$. Most importantly, the example supplies no map to three
positive lepton roots and does not evade the readout restrictions above.

A physical completion must therefore specify the carrier action,
generation embedding, admissible state domain, observable map and
dynamics on the same represented system. It must then determine which
parameters are predictions and which are measured inputs. Until those
steps are supplied, the framework supports conditional constructions
and consistency tests rather than a recovered particle spectrum or
a unified field theory. The long Companion retains the broader
cosmological and structural proposals with their individual qualifications.
